\begin{ex}
Cho dãy số liệu điểm thi: 7, 8, 9, 8, 10, 9, 8. Điểm số xuất hiện nhiều nhất là điểm mấy?
\shortans{8}
\loigiai{
Điểm 8 xuất hiện 3 lần.
}
\end{ex}

\begin{ex}
Tổng điểm của dãy số liệu trên là bao nhiêu?
\shortans{59}
\loigiai{
$7 + 8 + 9 + 8 + 10 + 9 + 8 = 59$.
}
\end{ex}

\begin{ex}
Biểu đồ cột số học sinh thích màu sắc: Đỏ: 25, Xanh: 30, Vàng: 20. Tổng số học sinh là bao nhiêu?
\shortans{75}
\loigiai{
$25 + 30 + 20 = 75$ bạn.
}
\end{ex}

\begin{ex}
Số học sinh thích màu xanh nhiều hơn màu vàng bao nhiêu bạn?
\shortans{10}
\loigiai{
$30 - 20 = 10$ bạn.
}
\end{ex}

\begin{ex}
Tung một đồng xu 50 lần, mặt Sấp xuất hiện 28 lần. Mặt Ngửa xuất hiện bao nhiêu lần?
\shortans{22}
\loigiai{
$50 - 28 = 22$ lần.
}
\end{ex}

\begin{ex}
Hiệu số lần giữa mặt Sấp và mặt Ngửa trong thí nghiệm trên là bao nhiêu?
\shortans{6}
\loigiai{
$28 - 22 = 6$ lần.
}
\end{ex}

\begin{ex}
Gieo xúc xắc 40 lần, mặt 6 chấm xuất hiện 8 lần. Tỉ số số lần xuất hiện mặt 6 chấm là một phần mấy? (Ghi mẫu số).
\shortans{5}
\loigiai{
$8 : 40 = 1/5$ (mẫu số là 5).
}
\end{ex}

\begin{ex}
Biểu đồ cột biểu diễn số cây trồng: Lớp 4A: 40 cây, Lớp 4B: 50 cây, Lớp 4C: 45 cây. Lớp trồng nhiều nhất trồng bao nhiêu cây?
\shortans{50}
\loigiai{
Lớp 4B: 50 cây.
}
\end{ex}

\begin{ex}
Tổng số cây cả 3 lớp trồng được là bao nhiêu?
\shortans{135}
\loigiai{
$40 + 50 + 45 = 135$ cây.
}
\end{ex}

\begin{ex}
Trung bình mỗi lớp trồng được bao nhiêu cây?
\shortans{45}
\loigiai{
$135 : 3 = 45$ cây.
}
\end{ex}

\begin{ex}
Số lượng sách đọc trong tuần của 4 bạn: 3, 5, 4, 8 cuốn. Bạn đọc nhiều nhất đọc bao nhiêu cuốn?
\shortans{8}
\loigiai{
Nhiều nhất là 8 cuốn.
}
\end{ex}

\begin{ex}
Tổng số sách cả 4 bạn đọc được là bao nhiêu cuốn?
\shortans{20}
\loigiai{
$3 + 5 + 4 + 8 = 20$ cuốn.
}
\end{ex}

\begin{ex}
Trung bình mỗi bạn đọc được bao nhiêu cuốn sách?
\shortans{5}
\loigiai{
$20 : 4 = 5$ cuốn.
}
\end{ex}

\begin{ex}
Tung đồng xu 60 lần, mặt Ngửa xuất hiện 32 lần. Mặt Sấp xuất hiện bao nhiêu lần?
\shortans{28}
\loigiai{
$60 - 32 = 28$ lần.
}
\end{ex}

\begin{ex}
Hai bạn chơi cờ vua 10 ván, bạn A thắng 6 ván, hòa 2 ván. Bạn B thắng mấy ván?
\shortans{2}
\loigiai{
$10 - (6 + 2) = 2$ ván.
}
\end{ex}

\begin{ex}
Gieo xúc xắc 20 lần, mặt chẵn xuất hiện 11 lần. Mặt lẻ xuất hiện bao nhiêu lần?
\shortans{9}
\loigiai{
$20 - 11 = 9$ lần.
}
\end{ex}

\begin{ex}
Trong biểu đồ cột, khoảng cách giữa các cột có bằng nhau không? (Có ghi 1, không ghi 0).
\shortans{1}
\loigiai{
Có (1).
}
\end{ex}

\begin{ex}
Số liệu ghi trên đỉnh mỗi cột biểu thị giá trị của cột đó đúng hay sai? (Đúng ghi 1, sai ghi 0).
\shortans{1}
\loigiai{
Đúng (1).
}
\end{ex}

\begin{ex}
Dãy số liệu gồm 5 số: 10, 15, 20, 25, 30. Số trung bình cộng là bao nhiêu?
\shortans{20}
\loigiai{
Số chính giữa là 20.
}
\end{ex}

\begin{ex}
Hiệu của số lớn nhất và số bé nhất trong dãy trên là bao nhiêu?
\shortans{20}
\loigiai{
$30 - 10 = 20$.
}
\end{ex}

\begin{ex}
Tung đồng xu 100 lần, mặt Sấp xuất hiện 55 lần. Mặt Ngửa ít hơn mặt Sấp bao nhiêu lần?
\shortans{10}
\loigiai{
Ngửa = 45 lần. Hiệu = $55 - 45 = 10$ lần.
}
\end{ex}

\begin{ex}
Một cửa hàng thống kê số quạt bán trong 4 ngày nắng nóng: 25, 30, 35, 30 chiếc. Cả 4 ngày bán được bao nhiêu chiếc?
\shortans{120}
\loigiai{
$25 + 30 + 35 + 30 = 120$ chiếc.
}
\end{ex}

\begin{ex}
Trung bình mỗi ngày cửa hàng bán được bao nhiêu chiếc quạt?
\shortans{30}
\loigiai{
$120 : 4 = 30$ chiếc.
}
\end{ex}

\begin{ex}
Một hộp có 5 bi xanh và 5 bi đỏ. Rút thăm 20 lần (có hoàn lại) được 12 lần bi đỏ. Số lần rút được bi xanh là bao nhiêu?
\shortans{8}
\loigiai{
$20 - 12 = 8$ lần.
}
\end{ex}

\begin{ex}
Tung xúc xắc 15 lần, số lần xuất hiện mặt 1 chấm nhiều nhất có thể là bao nhiêu?
\shortans{15}
\loigiai{
Nhiều nhất là 15 lần.
}
\end{ex}
